%
% 6.006 problem set 7 solutions template
%
\documentclass[12pt,twoside]{article}

\newcommand{\name}{}

\usepackage{amssymb}
\usepackage{amsmath}
\usepackage{graphicx}
\usepackage{latexsym}
\usepackage{times,url}
\usepackage{cprotect}
\usepackage{listings}
\usepackage{graphicx}
\usepackage[table]{xcolor}
\usepackage[letterpaper]{geometry}
\usepackage{tikz-qtree}
\usepackage{enumerate}

\newcommand{\profs}{Instructors: Erik Demaine, Jason Ku, and Justin Solomon}
\newcommand{\subj}{6.006}
\newcommand{\ttt}[1]{{\tt\small #1}}

\definecolor{dkgreen}{rgb}{0,0.6,0}
\definecolor{dkblue}{rgb}{.2,.2,1}
\definecolor{gray}{rgb}{0.5,0.5,0.5}
\definecolor{mauve}{rgb}{0.58,0,0.82}

\lstset{
  language=Python,
  aboveskip=1pc,
  belowskip=1pc,
  basicstyle={\footnotesize\ttfamily},
  numbers=left,
  showstringspaces=false,
  numberstyle={\tiny\color{gray}\ttfamily},
  keywordstyle={\color{dkblue}\ttfamily},
  commentstyle={\color{dkgreen}\ttfamily},
  stringstyle={\color{mauve}\ttfamily},
}

% \lstset{
%   language=Python,
%   aboveskip=1pc,
%   belowskip=1pc,
%   basicstyle={\bf\color{white}\ttfamily},
%   numbers=left,
%   showstringspaces=false,
%   numberstyle={\bf\small\color{lightgray}\ttfamily},
%   keywordstyle={\bf\color{cyan}\ttfamily},
%   commentstyle={\bf\color{green}\ttfamily},
%   stringstyle={\bf\color{mauve}\ttfamily},
% }

\tikzset{
  % every node/.style={minimum width=2em,draw,circle},
  % level 1/.style={sibling distance=2cm},
  level distance=1cm,
  edge from parent/.style=
  {draw,edge from parent path={(\tikzparentnode) -- (\tikzchildnode)}},
}
\usetikzlibrary{shapes}

\newif\ifHideSolutions
\newcommand{\solution}[1]{\color{dkgreen}\textbf{Solution: }#1\color{black}}
\newcommand{\commonmistakes}[1]{{\color{dkblue}\textbf{Common Mistakes: }#1}}
\newcommand{\rubric}[1]{\color{dkgreen}{\bf Rubric:} #1\color{black}}

% \HideSolutionsfalse
% \ifHideSolutions
%   \renewcommand{\solution}[1]{}
%   \renewcommand{\rubric}[1]{}
% \fi

\newlength{\toppush}
\setlength{\toppush}{2\headheight}
\addtolength{\toppush}{\headsep}

\newcommand{\htitle}[2]{\noindent\vspace*{-\toppush}\newline\parbox{\textwidth}
{\textit{Introduction to Algorithms: 6.006}\hfill\name\newline
Massachusetts Institute of Technology \hfill #2\newline
\profs\hfill #1 \\[-3.5ex]\newline
\mbox{}\hrulefill\mbox{}}\vspace*{1ex}\mbox{}\newline
\begin{center}{\Large\bf #1}\end{center}}

\newcommand{\handout}[2]{\thispagestyle{empty}
 \markboth{#1}{#1}
 \pagestyle{myheadings}\htitle{#1}{#2}}

\newcommand{\lecture}[3]{\thispagestyle{empty}
 \markboth{Lecture #1: #2}{Lecture #1: #2}
 \pagestyle{myheadings}\htitle{Lecture #1: #2}{#3}}

\newcommand{\htitlewithouttitle}[2]{\noindent\vspace*{-\toppush}\newline\parbox{6.5in}
{\textit{Introduction to Algorithms}\hfill#2\newline
Massachusetts Institute of Technology \hfill 6.006\newline
\profs\hfill Handout #1\vspace*{-.5ex}\newline
\mbox{}\hrulefill\mbox{}}\vspace*{1ex}\mbox{}\newline}

\newcommand{\handoutwithouttitle}[2]{\thispagestyle{empty}
 \markboth{Handout \protect\ref{#1}}{Handout \protect\ref{#1}}
 \pagestyle{myheadings}\htitlewithouttitle{\protect\ref{#1}}{#2}}

\newcommand{\exam}[2]{% parameters: exam name, date
 \thispagestyle{empty}
 \markboth{\hspace{1cm}\subj\ #1\hspace{1in}Name\hrulefill\ \ }%
          {\subj\ #1\hspace{1in}Name\hrulefill\ \ }
 \pagestyle{myheadings}\examtitle{#1}{#2}
 \renewcommand{\theproblem}{Problem \arabic{problemnum}}
}
\newcommand{\examsolutions}[3]{% parameters: handout, exam name, date
 \thispagestyle{empty}
 \markboth{Handout \protect\ref{#1}: #2}{Handout \protect\ref{#1}: #2}
% \pagestyle{myheadings}\htitle{\protect\ref{#1}}{#2}{#3}
 \pagestyle{myheadings}\examsolutionstitle{\protect\ref{#1}} {#2}{#3}
 \renewcommand{\theproblem}{Problem \arabic{problemnum}}
}
\newcommand{\examsolutionstitle}[3]{\noindent\vspace*{-\toppush}\newline\parbox{6.5in}
{\textit{Introduction to Algorithms}\hfill#3\newline
Massachusetts Institute of Technology \hfill 6.006\newline
%Singapore-MIT Alliance \hfill SMA5503\newline
\profs\hfill Handout #1\vspace*{-.5ex}\newline
\mbox{}\hrulefill\mbox{}}\vspace*{1ex}\mbox{}\newline
\begin{center}{\Large\bf #2}\end{center}}

\newcommand{\takehomeexam}[2]{% parameters: exam name, date
 \thispagestyle{empty}
 \markboth{\subj\ #1\hfill}{\subj\ #1\hfill}
 \pagestyle{myheadings}\examtitle{#1}{#2}
 \renewcommand{\theproblem}{Problem \arabic{problemnum}}
}

\makeatletter
\newcommand{\exambooklet}[2]{% parameters: exam name, date
 \thispagestyle{empty}
 \markboth{\subj\ #1}{\subj\ #1}
 \pagestyle{myheadings}\examtitle{#1}{#2}
 \renewcommand{\theproblem}{Problem \arabic{problemnum}}
 \renewcommand{\problem}{\newpage
 \item \let\@currentlabel=\theproblem
 \markboth{\subj\ #1, \theproblem}{\subj\ #1, \theproblem}}
}
\makeatother


\newcommand{\examtitle}[2]{\noindent\vspace*{-\toppush}\newline\parbox{6.5in}
{\textit{Introduction to Algorithms}\hfill#2\newline
Massachusetts Institute of Technology \hfill 6.006 Spring 2020\newline
%Singapore-MIT Alliance \hfill SMA5503\newline
\profs\hfill #1\vspace*{-.5ex}\newline
\mbox{}\hrulefill\mbox{}}\vspace*{1ex}\mbox{}\newline
\begin{center}{\Large\bf #1}\end{center}}

\newcommand{\grader}[1]{\hspace{1cm}\textsf{\textbf{#1}}\hspace{1cm}}

\newcommand{\points}[1]{[#1 points]\ }
\newcommand{\parts}[1]
{
  \ifnum#1=1
  (1 part)
  \else
  (#1 parts)
  \fi
  \ 
}

\newcommand{\bparts}{\begin{problemparts}}
\newcommand{\eparts}{\end{problemparts}}
\newcommand{\ppart}{\problempart}

%\newcommand{\lg} {lg\ }

\setlength{\oddsidemargin}{0pt}
\setlength{\evensidemargin}{0pt}
\setlength{\textwidth}{6.5in}
\setlength{\topmargin}{0in}
\setlength{\textheight}{8.5in}


\newcommand{\Spawn}{{\bf spawn} }
\newcommand{\Sync}{{\bf sync}}

\renewcommand{\cases}[1]{\left\{ \begin{array}{ll}#1\end{array}\right.}
\newcommand{\cif}[1]{\mbox{if $#1$}}
\newcommand{\cwhen}[1]{\mbox{when $#1$}}

\newcounter{problemnum}
\newcommand{\theproblem}{Problem \theproblemsetnum-\arabic{problemnum}}
\newenvironment{problems}{
        \begin{list}{{\bf \theproblem. \hspace*{0.5em}}}
        {\setlength{\leftmargin}{0em}
         \setlength{\rightmargin}{0em}
         \setlength{\labelwidth}{0em}
         \setlength{\labelsep}{0em}
         \usecounter{problemnum}}}{\end{list}}
\makeatletter
\newcommand{\problem}[1][{}]{\item \let\@currentlabel=\theproblem \textbf{#1}}
\makeatother

\newcounter{problempartnum}[problemnum]
\newenvironment{problemparts}{
        \begin{list}{{\bf (\alph{problempartnum})}}
        {\setlength{\leftmargin}{2.5em}
         \setlength{\rightmargin}{2.5em}
         \setlength{\labelsep}{0.5em}}}{\end{list}}
\newcommand{\problempart}{\addtocounter{problempartnum}{1}\item}

\newenvironment{truefalseproblemparts}{
        \begin{list}{{\bf (\alph{problempartnum})\ \ \ T\ \ F\hfil}}
        {\setlength{\leftmargin}{4.5em}
         \setlength{\rightmargin}{2.5em}
         \setlength{\labelsep}{0.5em}
         \setlength{\labelwidth}{4.5em}}}{\end{list}}

\newcounter{exercisenum}
\newcommand{\theexercise}{Exercise \theproblemsetnum-\arabic{exercisenum}}
\newenvironment{exercises}{
        \begin{list}{{\bf \theexercise. \hspace*{0.5em}}}
        {\setlength{\leftmargin}{0em}
         \setlength{\rightmargin}{0em}
         \setlength{\labelwidth}{0em}
         \setlength{\labelsep}{0em}
        \usecounter{exercisenum}}}{\end{list}}
\makeatletter
\newcommand{\exercise}{\item \let\@currentlabel=\theexercise}
\makeatother

\newcounter{exercisepartnum}[exercisenum]
%\newcommand{\problem}[1]{\medskip\mbox{}\newline\noindent{\bf Problem #1.}\hspace*{1em}}
%\newcommand{\exercise}[1]{\medskip\mbox{}\newline\noindent{\bf Exercise #1.}\hspace*{1em}}

\newenvironment{exerciseparts}{
        \begin{list}{{\bf (\alph{exercisepartnum})}}
        {\setlength{\leftmargin}{2.5em}
         \setlength{\rightmargin}{2.5em}
         \setlength{\labelsep}{0.5em}}}{\end{list}}
\newcommand{\exercisepart}{\addtocounter{exercisepartnum}{1}\item}


% Macros to make captions print with small type and 'Figure xx' in bold.
\makeatletter
\def\fnum@figure{{\bf Figure \thefigure}}
\def\fnum@table{{\bf Table \thetable}}
\let\@mycaption\caption
%\long\def\@mycaption#1[#2]#3{\addcontentsline{\csname
%  ext@#1\endcsname}{#1}{\protect\numberline{\csname 
%  the#1\endcsname}{\ignorespaces #2}}\par
%  \begingroup
%    \@parboxrestore
%    \small
%    \@makecaption{\csname fnum@#1\endcsname}{\ignorespaces #3}\par
%  \endgroup}
%\def\mycaption{\refstepcounter\@captype \@dblarg{\@mycaption\@captype}}
%\makeatother
\let\mycaption\caption
%\newcommand{\figcaption}[1]{\mycaption[]{#1}}

\newcounter{totalcaptions}
\newcounter{totalart}

\newcommand{\figcaption}[1]{\addtocounter{totalcaptions}{1}\caption[]{#1}}

% \psfigures determines what to do for figures:
%       0 means just leave vertical space
%       1 means put a vertical rule and the figure name
%       2 means insert the PostScript version of the figure
%       3 means put the figure name flush left or right
\newcommand{\psfigures}{0}
\newcommand{\spacefigures}{\renewcommand{\psfigures}{0}}
\newcommand{\rulefigures}{\renewcommand{\psfigures}{1}}
\newcommand{\macfigures}{\renewcommand{\psfigures}{2}}
\newcommand{\namefigures}{\renewcommand{\psfigures}{3}}

\newcommand{\figpart}[1]{{\bf (#1)}\nolinebreak[2]\relax}
\newcommand{\figparts}[2]{{\bf (#1)--(#2)}\nolinebreak[2]\relax}


\macfigures     % STATE

% When calling \figspace, make sure to leave a blank line afterward!!
% \widefigspace is for figures that are more than 28pc wide.
\newlength{\halffigspace} \newlength{\wholefigspace}
\newlength{\figruleheight} \newlength{\figgap}
\newcommand{\setfiglengths}{\ifnum\psfigures=1\setlength{\figruleheight}{\hruleheight}\setlength{\figgap}{1em}\else\setlength{\figruleheight}{0pt}\setlength{\figgap}{0em}\fi}
\newcommand{\figspace}[2]{\ifnum\psfigures=0\leavefigspace{#1}\else%
\setfiglengths%
\setlength{\wholefigspace}{#1}\setlength{\halffigspace}{.5\wholefigspace}%
\rule[-\halffigspace]{\figruleheight}{\wholefigspace}\hspace{\figgap}#2\fi}
\newlength{\widefigspacewidth}
% Make \widefigspace put the figure flush right on the text page.
\newcommand{\widefigspace}[2]{
\ifnum\psfigures=0\leavefigspace{#1}\else%
\setfiglengths%
\setlength{\widefigspacewidth}{28pc}%
\addtolength{\widefigspacewidth}{-\figruleheight}%
\setlength{\wholefigspace}{#1}\setlength{\halffigspace}{.5\wholefigspace}%
\makebox[\widefigspacewidth][r]{#2\hspace{\figgap}}\rule[-\halffigspace]{\figruleheight}{\wholefigspace}\fi}
\newcommand{\leavefigspace}[1]{\setlength{\wholefigspace}{#1}\setlength{\halffigspace}{.5\wholefigspace}\rule[-\halffigspace]{0em}{\wholefigspace}}

% Commands for including figures with macpsfig.
% To use these commands, documentstyle ``macpsfig'' must be specified.
\newlength{\macfigfill}
\makeatother
\newlength{\bbx}
\newlength{\bby}
\newcommand{\macfigure}[5]{\addtocounter{totalart}{1}
\ifnum\psfigures=2%
\setlength{\bbx}{#2}\addtolength{\bbx}{#4}%
\setlength{\bby}{#3}\addtolength{\bby}{#5}%
\begin{flushleft}
\ifdim#4>28pc\setlength{\macfigfill}{#4}\addtolength{\macfigfill}{-28pc}\hspace*{-\macfigfill}\fi%
\mbox{\psfig{figure=./#1.ps,%
bbllx=#2,bblly=#3,bburx=\bbx,bbury=\bby}}
\end{flushleft}%
\else\ifdim#4>28pc\widefigspace{#5}{#1}\else\figspace{#5}{#1}\fi\fi}
\makeatletter

\newlength{\savearraycolsep}
\newcommand{\narrowarray}[1]{\setlength{\savearraycolsep}{\arraycolsep}\setlength{\arraycolsep}{#1\arraycolsep}}
\newcommand{\normalarray}{\setlength{\arraycolsep}{\savearraycolsep}}

\newcommand{\hint}{{\bf Hint:\ }}

% Macros from /th/u/clr/mac.tex

\newcommand{\set}[1]{\left\{ #1 \right\}}
\newcommand{\abs}[1]{\left| #1\right|}
\newcommand{\card}[1]{\left| #1\right|}
\newcommand{\floor}[1]{\left\lfloor #1 \right\rfloor}
\newcommand{\ceil}[1]{\left\lceil #1 \right\rceil}
\newcommand{\ang}[1]{\ifmmode{\left\langle #1 \right\rangle}
   \else{$\left\langle${#1}$\right\rangle$}\fi}
        % the \if allows use outside mathmode,
        % but will swallow following space there!
\newcommand{\paren}[1]{\left( #1 \right)}
\newcommand{\bracket}[1]{\left[ #1 \right]}
\newcommand{\prob}[1]{\Pr\left\{ #1 \right\}}
\newcommand{\Var}{\mathop{\rm Var}\nolimits}
\newcommand{\expect}[1]{{\rm E}\left[ #1 \right]}
\newcommand{\expectsq}[1]{{\rm E}^2\left[ #1 \right]}
\newcommand{\variance}[1]{{\rm Var}\left[ #1 \right]}
\renewcommand{\choose}[2]{{{#1}\atopwithdelims(){#2}}}
\def\pmod#1{\allowbreak\mkern12mu({\rm mod}\,\,#1)}
\newcommand{\matx}[2]{\left(\begin{array}{*{#1}{c}}#2\end{array}\right)}
\newcommand{\Adj}{\mathop{\rm Adj}\nolimits}

\newtheorem{theorem}{Theorem}
\newtheorem{lemma}[theorem]{Lemma}
\newtheorem{corollary}[theorem]{Corollary}
\newtheorem{xample}{Example}
\newtheorem{definition}{Definition}
\newenvironment{example}{\begin{xample}\rm}{\end{xample}}
\newcommand{\proof}{\noindent{\em Proof.}\hspace{1em}}
\def\squarebox#1{\hbox to #1{\hfill\vbox to #1{\vfill}}}
\newcommand{\qedbox}{\vbox{\hrule\hbox{\vrule\squarebox{.667em}\vrule}\hrule}}
\newcommand{\qed}{\nopagebreak\mbox{}\hfill\qedbox\smallskip}
\newcommand{\eqnref}[1]{(\protect\ref{#1})}

%%\newcommand{\twodots}{\mathinner{\ldotp\ldotp}}
\newcommand{\transpose}{^{\mbox{\scriptsize \sf T}}}
\newcommand{\amortized}[1]{\widehat{#1}}

\newcommand{\punt}[1]{}

%%% command for putting definitions into boldface
% New style for defined terms, as of 2/23/88, redefined by THC.
\newcommand{\defn}[1]{{\boldmath\textit{\textbf{#1}}}}
\newcommand{\defi}[1]{{\textit{\textbf{#1\/}}}}

\newcommand{\red}{\leq_{\rm P}}
\newcommand{\lang}[1]{%
\ifmmode\mathord{\mathcode`-="702D\rm#1\mathcode`\-="2200}\else{\rm#1}\fi}

%\newcommand{\ckt}[1]{\ifmmode\mathord{\mathcode`-="702D\sc #1\mathcode`\-="2200}\else$\mathord{\mathcode`-="702D\sc #1\mathcode`\-="2200}$\fi}
\newcommand{\ckt}[1]{\ifmmode \sc #1\else$\sc #1$\fi}

%% Margin notes - use \notesfalse to turn off notes.
\setlength{\marginparwidth}{0.6in}
\reversemarginpar
\newif\ifnotes
\notestrue
\newcommand{\longnote}[1]{
  \ifnotes
    {\medskip\noindent Note: \marginpar[\hfill$\Longrightarrow$]
      {$\Longleftarrow$}{#1}\medskip}
  \fi}
\newcommand{\note}[1]{
  \ifnotes
    {\marginpar{\tiny \raggedright{#1}}}
  \fi}


\newcommand{\reals}{\mathbbm{R}}
\newcommand{\integers}{\mathbbm{Z}}
\newcommand{\naturals}{\mathbbm{N}}
\newcommand{\rationals}{\mathbbm{Q}}
\newcommand{\complex}{\mathbbm{C}}

\newcommand{\oldreals}{{\bf R}}
\newcommand{\oldintegers}{{\bf Z}}
\newcommand{\oldnaturals}{{\bf N}}
\newcommand{\oldrationals}{{\bf Q}}
\newcommand{\oldcomplex}{{\bf C}}

\newcommand{\w}{\omega}                 %% for fft chapter

\newenvironment{closeitemize}{\begin{list}
{$\bullet$}
{\setlength{\itemsep}{-0.2\baselineskip}
\setlength{\topsep}{0.2\baselineskip}
\setlength{\parskip}{0pt}}}
{\end{list}}

% These are necessary within a {problems} environment in order to restore
% the default separation between bullets and items.
\newenvironment{normalitemize}{\setlength{\labelsep}{0.5em}\begin{itemize}}
                              {\end{itemize}}
\newenvironment{normalenumerate}{\setlength{\labelsep}{0.5em}\begin{enumerate}}
                                {\end{enumerate}}

%\def\eqref#1{Equation~(\ref{eq:#1})}
%\newcommand{\eqref}[1]{Equation (\ref{eq:#1})}
\newcommand{\eqreftwo}[2]{Equations (\ref{eq:#1}) and~(\ref{eq:#2})}
\newcommand{\ineqref}[1]{Inequality~(\ref{ineq:#1})}
\newcommand{\ineqreftwo}[2]{Inequalities (\ref{ineq:#1}) and~(\ref{ineq:#2})}

\newcommand{\figref}[1]{Figure~\ref{fig:#1}}
\newcommand{\figreftwo}[2]{Figures \ref{fig:#1} and~\ref{fig:#2}}

\newcommand{\liref}[1]{line~\ref{li:#1}}
\newcommand{\Liref}[1]{Line~\ref{li:#1}}
\newcommand{\lirefs}[2]{lines \ref{li:#1}--\ref{li:#2}}
\newcommand{\Lirefs}[2]{Lines \ref{li:#1}--\ref{li:#2}}
\newcommand{\lireftwo}[2]{lines \ref{li:#1} and~\ref{li:#2}}
\newcommand{\lirefthree}[3]{lines \ref{li:#1}, \ref{li:#2}, and~\ref{li:#3}}

\newcommand{\lemlabel}[1]{\label{lem:#1}}
\newcommand{\lemref}[1]{Lemma~\ref{lem:#1}} 

\newcommand{\exref}[1]{Exercise~\ref{ex:#1}}

\newcommand{\handref}[1]{Handout~\ref{#1}}

\newcommand{\defref}[1]{Definition~\ref{def:#1}}

% (1997.8.16: Victor Luchangco)
% Modified \hlabel to only get date and to use handouts counter for number.
%   New \handout and \handoutwithouttitle commands in newmac.tex use this.
%   The date is referenced by <label>-date.
%   (Retained old definition as \hlabelold.)
%   Defined \hforcelabel to use an argument instead of the handouts counter.

\newcounter{handouts}
\setcounter{handouts}{0}

\newcommand{\hlabel}[2]{%
\stepcounter{handouts}
{\edef\next{\write\@auxout{\string\newlabel{#1}{{\arabic{handouts}}{0}}}}\next}
\write\@auxout{\string\newlabel{#1-date}{{#2}{0}}}
}

\newcommand{\hforcelabel}[3]{%          Does not step handouts counter.
\write\@auxout{\string\newlabel{#1}{{#2}{0}}}
\write\@auxout{\string\newlabel{#1-date}{{#3}{0}}}}


% less ugly underscore
% --juang, 2008 oct 05
\renewcommand{\_}{\vrule height 0 pt depth 0.4 pt width 0.5 em \,}

% multiline framed box (will always extend to the far right edge; for a short single line, use \fbox directly)
% --zabel, fall 2018
\newcommand\framepar[1]{\fbox{\begin{minipage}{\linewidth}#1\end{minipage}}}


\newcommand{\theproblemsetnum}{7}

\title{6.006 Problem Set 7}

\begin{document}

\handout{Problem Set \theproblemsetnum}

\setlength{\parindent}{0pt}
\medskip\hrulefill\medskip

{\bf Name:} Gabriel Chiong

\medskip

{\bf Collaborators:} None

\medskip\hrulefill

%%%%%%%%%%%%%%%%%%%%%%%%%%%%%%%%%%%%%%%%%%%%%%%%%%%%%
% See below for common and useful latex constructs. %
%%%%%%%%%%%%%%%%%%%%%%%%%%%%%%%%%%%%%%%%%%%%%%%%%%%%%

% Some useful commands:
%$f(x) = \Theta(x)$
%$T(x, y) \leq \log(x) + 2^y + \binom{2n}{n}$
% {\tt code\_function}


% You can create unnumbered lists as follows:
%\begin{itemize}
%    \item First item in a list
%        \begin{itemize}
%            \item First item in a list
%                \begin{itemize}
%                    \item First item in a list
%                    \item Second item in a list
%                \end{itemize}
%            \item Second item in a list
%        \end{itemize}
%    \item Second item in a list
%\end{itemize}

% You can create numbered lists as follows:
%\begin{enumerate}
%    \item First item in a list
%    \item Second item in a list
%    \item Third item in a list
%\end{enumerate}

% You can write aligned equations as follows:
%\begin{align}
%    \begin{split}
%        (x+y)^3 &= (x+y)^2(x+y) \\
%                &= (x^2+2xy+y^2)(x+y) \\
%                &= (x^3+2x^2y+xy^2) + (x^2y+2xy^2+y^3) \\
%                &= x^3+3x^2y+3xy^2+y^3
%    \end{split}
%\end{align}

% You can create grids/matrices as follows:
%\begin{align}
%    A =
%    \begin{bmatrix}
%        A_{11} & A_{21} \\
%        A_{21} & A_{22}
%    \end{bmatrix}
%\end{align}

% You can include images and PDFs as follows:
% \includegraphics[width=0.5\textwidth]{img.jpg}

\begin{problems}

\problem  % Problem 1

Subproblems:
\begin{itemize}
    \item Let $x(i)$ denote the maximum delegates that Torr can win in state $i$ for all $i \in \{1,\ldots,n+1\}$.
\end{itemize}

Relate:
\begin{itemize}
    \item Torr can either choose to campaign or not on any day.
    \item If she chooses to campaign on day $i$, then $x_c(i)=d_i+z_{i+1}+z_{i+2}+x(i+3)$.
    \item If she chooses to do nothing on day $i$, then $x_n(i)=z_i+x(i+1)$.
    \item Therefore the overall relation can be written as $x(i)=\max\{x_c(i),x_n(i)\}$.
\end{itemize}

Topological order:
\begin{itemize}
    \item $x(i)$ only depends on subproblems with strictly larger $i$, therefore the relation graph is a DAG.
\end{itemize}

Base:
\begin{itemize}
    \item For $i>n$, we have $x(i)=0$.
    \item For $i=n$, we have $x(i)=d_n$.
    \item For $i=n-1$, we have $x(i)=\max\{d_i+z_n, z_i+d_n\}$.
\end{itemize}

Original:
\begin{itemize}
    \item Calculate $x(1)$ and check if $x(1)>\lfloor D/2 \rfloor + 1$.
\end{itemize}

Time:
\begin{itemize}
    \item There are $n+1$ subproblems including the base case, and constant $O(1)$ work done per sub-problem. Therefore, the overall running time is $O(n)$.
\end{itemize}

\newpage
\problem  % Problem 2
We first sort the tigers by age using an $O(n\log n)$ algorithm like merge sort. Next, sort the cages by distance in increasing order in $O(n^2\log n)$ time, again using merge sort.

Subproblems:
\begin{itemize}
    \item Define $x(i,j)$ be the minimum total discomfort for tigers \verb|T[i:]| in cages \verb|C[j:]| for all $i \in \{0,\ldots,n\}$ and $j \in \{0,\ldots,n^2\}$.
\end{itemize}

Relate:
\begin{itemize}
    \item Let the discomfort $d(i,j)$ of tiger $i$ in cage $j$ be $s_i-c_j$ if $s_i>c_j$ and 0 otherwise.
    \item We can choose to match every $i,j$ or not, therefore we have $x(i,j)=\min\{d(i,j)+x(i+1,j+1),x(i,j+1)\}$.
\end{itemize}

Topological order:
\begin{itemize}
    \item $x(i,j)$ at any stage relies on strictly larger subproblem $j$s, therefore the relation graph is a DAG.
\end{itemize}

Base:
\begin{itemize}
    \item $x(n,j)=0$, since there are no more tigers are left to match, there can be no more additional discomfort.
    \item $x(i,n^2)=\infty$, to penalize the scenario of not matching all tigers to cages.
\end{itemize}

Original:
\begin{itemize}
    \item $x(0,0)$ to calculate the discomfort considering all tigers and all cages.
    \item We can store parent pointers to reconstruct the optimal assignment.
\end{itemize}

Time:
\begin{itemize}
    \item There are $(n+1)(n^2+1)=O(n^3)$ subproblems.
    \item Each subproblem takes a constant $O(1)$ amount of work.
    \item Therefore the overall algorithm runs in $O(n^3)$ time as required.
\end{itemize}

\newpage
\problem  % Problem 3

Subproblems:
\begin{itemize}
    \item We define $i \in \{0, 1\}$ as an indicator of the parity of the number of even or odd paths.
    \item Let $x(v,i)$ be the number of even paths from $(s,v)$ if $i=0$ and the number of odd paths is $i=1$ for all $v \in V$.
    \item Therefore, there are effectively two vertices for every vertex in the original graph $G$.
\end{itemize}

Relate:
\begin{itemize}
    \item Let $\chi(k)$ be a function with returns 1 if $k$ is even, and 0 otherwise.
    \item Then we can define $x(v,i)=\sum \{x(u, \chi(w(u,v)+i)):u \in Adj^-(v)\}$.
    \item $u$ is the set of all incoming vertices to $v$.
\end{itemize}

Topological order:
\begin{itemize}
    \item $x(v,i)$ depends only on $x(u,j)$ where $u$ appears before $v$ int he topological order of $G$.
\end{itemize}

Base:
\begin{itemize}
    \item $x(s,0)=1$, since $0$ is an even number.
    \item $x(s,1)=0$, for the case where there are no odd length paths from $(s,s)$.
    \item $x(v,0)=x(v,1)=0$ for all $v: Adj^-(v)=\emptyset$, where there are no incoming vertices to $v$.
\end{itemize}

Original:
\begin{itemize}
    \item $x(t,1)$ calculates the odd paths in reverse topological order.
\end{itemize}

Time:
\begin{itemize}
    \item There are $2|V|$ subproblems (a vertex for each odd/even pair).
    \item Each subproblem requires iterating over all incoming vertices, which is $O(\deg^-(v))$.
    \item Therefore, the overall runtime is $O(2 \sum_{v \in V} \deg^-(v))=O(|V|+|E|)$, which is linear in the input graph size.
\end{itemize}

\newpage
\problem  % Problem 4
As the game progresses, the slices eaten by round are cyclically consecutive (since the slices are eaten in $\alpha_i$ to $\alpha_i+\pi$). Therefore our subproblems can be consecutive cyclic subarrays. Since each sibling wishes to maximize tastiness, our subproblems require an indicator for which sibiling is currently making the choice.

Let $v(i,j)$ be the tastiness of the $j$ slices counter-clockwise from angle $\alpha_i$, specifically $v(i,j)=\sum_{k=0}^{j-1}t_{((i+k)\pmod{2n})}$, where $i\in \{0,\ldots,2n-1\}$ and $j\in \{0,\ldots,n\}$. There are $O(n^2)$ ways to construct a $v(i,j)$, and each can be computed in $O(n)$ time for an overall time of $O(n^3)$.

Subproblems:
\begin{itemize}
    \item Let $x(i,j,p)$ be the maximum tastiness Liza can achieve with the $j$ slices counter-clockwise from $\alpha_i$ remaining, when either Liza is the chooser ($p=1$), or Lie is the chooser ($p=2)$. 
    \item Also let $i,j,p$ be defined as $i \in \{0,\ldots,2n-1\}$, $j \in \{0,\ldots,n\}$, and $p\in \{1,2\}$.
\end{itemize}

Relate:
\begin{itemize}
    \item Chooser can choose any proper angle and then choose a side. Therefore, we consider all possibilities.
    \item Angle $\alpha_{i+k}$ is proper for any $k \in \{1,\ldots,j-1\}$.
    \item Chooser eats either $k$ slices between $\alpha_i$ and $\alpha_{i+k}$ or $j-k$ slices between $\alpha_{i+k}$ and $\alpha_{i+j}$.
    \item Neither gains or loses tastiness from the choice of the other.
    \item For Liza, $x(i,j,1)=\max\{\max\{v(i,k)+x(i+k,j-k,2), v(i+k,j-k)+x(i,k,2)\} : k \in \{1,\ldots,j-1\}\}$.
    \item For Lie, $x(i,j,2)=\min\{\min\{x(i+k,j-k,1),x(i,k,1)\} : k \in \{1,\ldots,j-1\}\}$.
\end{itemize}

Topological order:
\begin{itemize}
    \item $x(i,j,p)$ only depends on subproblems with strictly smaller $j$, therefore the relation graph is a DAG.
\end{itemize}

Base:
\begin{itemize}
    \item Liza eats the last slice, $x(i,1,1)=t_i$, for all $i \in \{1,\ldots,2n\}$.
    \item Lie eats the last slice, $x(i,1,2)=0$, for all $i \in \{1,\ldots,2n\}$.
\end{itemize}

Original
\begin{itemize}
    \item Liza starts the game, and the maximum tastiness on a half and letting Lie choose on the other half.
    \item $\max\{x(i,n,2)+v(((i+n)\pmod {2n}),n):i \in \{0,\ldots,2n-1\}\}$
\end{itemize}

Time:
\begin{itemize}
    \item There are $2(2n)(n+1)=O(n^2)$ subproblems in total, each requiring $O(n)$ work per subproblem. The original also requires $O(n)$ work, so the overall running time is $O(n^3)$.
\end{itemize}

\newpage
\problem  % Problem 5

Subproblems:
\begin{itemize}
    \item Let $x(j)$ be the longest decreasing subsequence of prices that doesn't skip days in \verb|P[j:| that includes price \verb|P[j]|, for all $j \in \{0,\ldots,nk-1\}$.
\end{itemize}

Relate:
\begin{itemize}
    \item Next stock price information in the sequence is at a later time in the same day or the next, so we iterate over all possibilities.
    \item At price index $j$, there are $(k-1)-(j \mod k)$ prices left in the same day, then there are $k$ prices the following day (if it exists).
    \item The last next price index is $f(j)=\min\{j+(k-1)-(j \mod k)+k,nk-1\}$.
    \item This results in the relation $x(j)=1+\max\{x(d):d\in\{j+1,\ldots,f(j)\} \text{ and } P[j]>P[d]\}\cup \{0\}$.
\end{itemize}

Topological order:
\begin{itemize}
    \item $x(j)$ only depends on subproblems with strictly larger $j$, so relation graph is a DAG.
\end{itemize}

Base:
\begin{itemize}
    \item When only one item remains, $x(nk-1)=1$.
\end{itemize}

Original:
\begin{itemize}
    \item The shorting value is $\max\{x(j): j\in \{0,\ldots,nk-1\}\}$.
\end{itemize}

Time:
\begin{itemize}
    \item There are $nk$ subproblems, and $O(k)$ work per subproblem. 
    \item The work for original is $O(nk)$, but is dominated by the work for the subproblems. 
    \item Therefore, the overall running time is $O(nk^2)$.
\end{itemize}

\begin{problemparts}
\problempart % Problem 5a
\problempart Submit your implementation to {\small\url{alg.mit.edu}}.
\end{problemparts}

\end{problems}

\end{document}
